% requer booktabs. Descritivo (BUILD09).
\begin{table}[htbp]\centering
\caption{Size of the eligible set: \emph{static} count of childcare centres within 2\,km of the postcode, paper universe, total and by year of application}
\label{tab:N2}
\small
\setlength{\tabcolsep}{5pt}
\begin{tabular}{lrrrrrrr}
\toprule
Group & p10 & p25 & Median & p75 & p90 & Mean & N \\
\midrule
Total (2006--2018) & 13 & 21 & 32 & 48 & 66 & 36.0 & 1,524,560 \\
\midrule
2010 & 13 & 21 & 33 & 49 & 67 & 36.7 & 81,004 \\
2011 & 13 & 21 & 33 & 48 & 66 & 36.1 & 124,240 \\
2012 & 13 & 21 & 33 & 49 & 68 & 36.6 & 220,922 \\
2013 & 13 & 21 & 33 & 48 & 67 & 36.2 & 191,149 \\
2014 & 13 & 21 & 32 & 48 & 66 & 35.9 & 174,676 \\
2015 & 13 & 21 & 32 & 48 & 66 & 35.9 & 179,869 \\
2016 & 13 & 21 & 32 & 47 & 65 & 35.6 & 172,341 \\
2017 & 12 & 21 & 32 & 47 & 65 & 35.5 & 163,922 \\
2018 & 13 & 20 & 31 & 47 & 65 & 35.2 & 165,331 \\
\bottomrule
\end{tabular}
\\[3pt]{\footnotesize \textit{Note:} Universe: all first applications (B\_novo, 1,524,560 children). The eligible set is the count of childcare centres within 2\,km of the child's postcode, from the postcode-to-centre distance file. This is the \emph{static} count: it includes every centre in the register regardless of whether it had opened at the application date, so year-to-year differences reflect the geography of where applicants live, not growth of the network. The preferred measure conditions on centres already open at the application date (a median of 18, rising from 12 in 2010 to 29 in 2018); it is reported, alongside this static count, in the dynamic eligible-set table.}
\end{table}
